\chapter{Байесовский вывод, априорные распределения и выбор моделей}
\label{ch:bayes}

Частотный подход трактует параметр $\theta$ как фиксированную константу, а случайность~--- в данных.
Байесовский подход объявляет $\theta$ случайной величиной с априорным распределением $p(\theta)$, обновляемым по наблюдениям до апостериорного $p(\theta\mid X)$ \citep{gelman2013,wasserman2004}.
Это даёт полное распределение неопределённости, достоверные интервалы, апостериорное прогнозное распределение для новых наблюдений и естественный механизм выбора моделей через маргинальное правдоподобие.

В гл.~\ref{ch:hypothesis} мы трактовали p-value и доверительные интервалы как свойства процедуры при фиксированном $\theta$. В гл.~\ref{ch:regression} гребневая и lasso-регуляризация уже выступали как штрафы к правдоподобию.
Байесовский вывод делает эту связь явной. Априор $p(\theta)$ и правдоподобие $p(X\mid\theta)$ объединяются в одну плотность, а неопределённость после данных выражается распределением, а не единственной точкой (МП или МНК).
Для последовательных моделей (гл.~\ref{ch:causality}, HMM) EM и Баума--Уэлча~--- частотные аналоги итеративного уточнения. MCMC и VI дают байесовскую альтернативу, когда апостериор не вычисляется аналитически.

\section{Формула Байеса и вывод}

Для параметра $\theta$ и данных $X=(X_1,\ldots,X_n)$: \[ p(\theta\mid X)=\frac{p(X\mid\theta)\,p(\theta)}{p(X)}, \qquad p(X)=\int p(X\mid\theta)\,p(\theta)\,\mathrm{d}\theta.
\] Формула Байеса~--- тот же закон условной вероятности, что лежит в основе цепи Маркова и HMM (гл.~\ref{ch:causality}). Апостериор пропорционален произведению правдоподобия на априор при фиксированных данных.
Логарифм апостериорной плотности: \[ \log p(\theta\mid X)=\log p(X\mid\theta)+\log p(\theta)+\mathrm{const}.
\] Нормировочная константа $p(X)$~--- маргинальное правдоподобие. Она не нужна для MAP и MCMC, но критична при сравнении моделей.

\begin{definition}
MAP-оценка $\hat\theta_{\mathrm{MAP}}=\arg\max_\theta p(\theta\mid X)$ совпадает с регуляризованным МП. Априорное $\log p(\theta)$ играет роль штрафа (аналогично ridge/lasso в гл.~\ref{ch:regression}).
\end{definition}

\begin{definition}
Семейство априоров сопряжено правдоподобию, если апостериор принадлежит тому же семейству, что и априор.
\end{definition}

\section{Пример: монетка и бета-априор}

В примере с монетой зафиксированы десять подбрасываний. Семь орлов и три решки. Модель: $X_i\mid p\sim \mathrm{Bern}(p)$ i.i.d.
Это та же схема успех/неуспех, что в биномиальных тестах (гл.~\ref{ch:hypothesis}), но параметр $p$ трактуется как случайная величина, а не как фиксированная частота генерации.

Априор $p\sim \mathrm{Beta}(\alpha,\beta)$.
Правдоподобие $p(X\mid p)=p^{\sum X_i}(1-p)^{n-\sum X_i}$.
Апостериор.\[ p\mid X \sim \mathrm{Beta}\!\left(\alpha+\sum X_i,\;\beta+n-\sum X_i\right).
\] MAP (при $\alpha,\beta>1$): $\hat p=\frac{\alpha+\sum X_i-1}{\alpha+\beta+n-2}$.
Апостериорное среднее $\mathbb{E}(p\mid X)=\frac{\alpha+\sum X_i}{\alpha+\beta+n}$.
При $\mathrm{Beta}(1,1)$ и $k=7$ орлов: $p\mid X\sim\mathrm{Beta}(8,4)$, $\mathbb{E}(p\mid X)=0{,}667$.
На рис.~\ref{fig:beta-posterior}~--- апостериорная плотность при симметричном априоре $\mathrm{Beta}(1,1)$ и данных 7H/3T.

\begin{figure}[htbp]
    \centering
    \includegraphics[width=\FigWidth]{\FigDir/chapter-8/beta.png}
    \caption{Апостериорное распределение $p$ (Beta).}
    \label{fig:beta-posterior}
\end{figure}

При пяти орлах подряд МП даёт $\hat p=0$ для вероятности решки. Байесовский вывод с умеренным априором (например, $\mathrm{Beta}(2,2)$) сохраняет ненулевую вероятность хвоста~--- это иллюстрация роли априора при малых $n$.
Достоверный интервал для $p$ при $\mathrm{Beta}(8,4)$ интерпретируется как вероятность параметра в интервале равна 95\% после наблюдения данных~--- в отличие от частотного ДИ (гл.~\ref{ch:basics},~\ref{ch:hypothesis}).
При 7H/3T и $\mathrm{Beta}(1,1)$ апостериор $\mathrm{Beta}(8,4)$. Апостериорное среднее $2/3$, HPD смещён относительно интервала с равными хвостами. Прогнозное $P(\text{next H})=8/12$.
Та же цепочка априор $\to$ данные $\to$ апостериор $\to$ прогноз повторяется в NIG и Gamma--Poisson.

\section{Сопряжённые семейства}

Нормальная модель: Normal--Inverse-Gamma. Данные $X_i\mid\mu,\sigma^2\sim\mathcal{N}(\mu,\sigma^2)$ i.i.d.
Сопряжённый априор на $(\mu,\sigma^2)$: \[ \mu\mid\sigma^2 \sim \mathcal{N}(\mu_0,\,\sigma^2/\kappa_0), \qquad \sigma^2 \sim \mathrm{Inv\text{-}Gamma}(\alpha_0,\,\beta_0).
\] Апостериор после $n$ наблюдений с $\bar x$, $s^2$: \[ \mu\mid\sigma^2,X \sim \mathcal{N}\!\left(\mu_n,\,\frac{\sigma^2}{\kappa_n}\right), \quad \sigma^2\mid X \sim \mathrm{Inv\text{-}Gamma}(\alpha_n,\beta_n), \] где $\kappa_n=\kappa_0+n$, $\mu_n=\frac{\kappa_0\mu_0+n\bar x}{\kappa_n}$, $\alpha_n=\alpha_0+n/2$, $\beta_n=\beta_0+\frac12\sum(x_i-\bar x)^2+\frac{\kappa_0 n(\bar x-\mu_0)^2}{2\kappa_n}$.
Маргинальный апостериор $\mu\mid X$~--- распределение Стьюдента. Достоверный интервал для $\mu$ не совпадает с частотным $t$-интервалом при том же $\alpha$, хотя при больших $n$ близок.
NIG-модель~--- байесовский аналог линейной регрессии с неизвестной дисперсией. Априор на $(\mu,\sigma^2)$ заменяет фиксированное $\sigma^2$ в классическом МНК и даёт $t$-подобные хвосты для $\mu\mid X$.

Пуассон и гамма. Счётчики $N_i\mid\lambda\sim\mathrm{Poisson}(\lambda)$: \[ p(N\mid\lambda)\propto \lambda^{\sum N_i} e^{-\lambda n}.
\] Сопряжённый априор: $\lambda\sim\mathrm{Gamma}(\alpha,\beta)$ (параметризация $\mathrm{rate}=\beta$): \[ \lambda\mid N \sim \mathrm{Gamma}\!\left(\alpha+\sum N_i,\;\beta+n\right).
\] Апостериорное среднее: $\mathbb{E}(\lambda\mid N)=\frac{\alpha+\sum N_i}{\beta+n}$.

Апостериорное прогнозное распределение (прогноз нового наблюдения) интегрирует правдоподобие по апостериору \[ p(X_{\mathrm{new}}\mid X)=\int p(X_{\mathrm{new}}\mid\theta)\,p(\theta\mid X)\,\mathrm{d}\theta.
\] Прогнозное распределение шире, чем подстановка $\hat\theta$. Оно учитывает неопределённость параметра~--- аналог доверительного прогнозного интервала в регрессии, но в чисто байесовской форме.
Beta--Binomial: при $p\mid X\sim\mathrm{Beta}(a,b)$ \[ P(X_{\mathrm{new}}=1\mid X)=\frac{a}{a+b}.
\] Нормальное--Стьюдента: при NIG-апостериоре предсказание для одного $X_{\mathrm{new}}$~--- $t$-распределение с тяжёлыми хвостами (учёт неопределённости $\sigma^2$).
Gamma--Poisson: предсказание~--- отрицательное биномиальное распределение для счётчиков.
Сопряжённость удобна аналитически, но не обязательна. Для произвольного априора используют MCMC, вариационный вывод (VI) или сеточные аппроксимации.
При большом $n$ апостериор определяется правдоподобием. Слабый априор подавляется данными. VI минимизирует $KL(q(\theta)\|p(\theta\mid X))$ внутри семейства $q$ (например, факторизованное гауссово). Быстрее MCMC, но смещение при узком $q$.
Для Beta--Binomial новый орёл учитывает неопределённость $p$, а не только $\bar x$.
В гл.~\ref{ch:basics} частотный интервал для $p$ описывает процедуру при фиксированном $p$. Прогнозное $P(X_{\mathrm{new}}=1\mid X)$~--- распределение будущего наблюдения при текущих данных.

\begin{center}
\footnotesize
\begin{tabular}{lllll}
\toprule
Правдоподобие & Параметр & Априор & Апостериор (параметры) \\
\midrule
Бернулли / биномиальное & $p$ & $\mathrm{Beta}(\alpha,\beta)$ &
$\mathrm{Beta}(\alpha{+}k,\,\beta{+}n{-}k)$, $k=\sum x_i$ \\
Нормальное, $\sigma^2$ известна & $\mu$ & $\mathcal{N}(\mu_0,\sigma_0^2)$ &
$\mathcal{N}(\mu_n,\sigma_n^2)$, $\mu_n=\frac{\sigma_0^{-2}\mu_0+n\sigma^{-2}\bar x}{\sigma_0^{-2}+n\sigma^{-2}}$ \\
Нормальное, $\mu$ известно & $\sigma^2$ & $\mathrm{Inv\text{-}}\Gamma(\alpha_0,\beta_0)$ &
$\mathrm{Inv\text{-}}\Gamma(\alpha_n,\beta_n)$ \\
Нормальное i.i.d. & $(\mu,\sigma^2)$ & NIG &
$\mu\mid\sigma^2,X\sim\mathcal{N}(\mu_n,\sigma^2/\kappa_n)$,
$\sigma^2\mid X\sim\mathrm{Inv\text{-}}\Gamma(\alpha_n,\beta_n)$ \\
Пуассон & $\lambda$ & $\mathrm{Gamma}(\alpha,\beta)$ (rate) &
$\mathrm{Gamma}(\alpha{+}{\sum N_i},\,\beta{+}n)$ \\
Мультиномиальное & $\mathbf{p}$ & $\mathrm{Dir}(\boldsymbol{\alpha})$ &
$\mathrm{Dir}(\boldsymbol{\alpha}+\mathbf{n})$ \\
\bottomrule
\end{tabular}
\end{center}

\section{Назначение априорных распределений}

Информативный априор отражает экспертные знания или прошлые эксперименты (например, $\mathrm{Beta}(2,2)$~--- умеренное смещение к $p=0{,}5$).
Неинформативный / слабоинформативный.~--- широкая поддержка ($\mathrm{Beta}(1,1)$, нормальный с большой дисперсией, равномерный на $[-50,50]$ для температуры).
Выбор между ними~--- вопрос контекста. В промышленном контроле качества разумен информативный априор от прошлых партий. В разведочном анализе~--- слабый, чтобы данные доминировали.
Распределение Джеффриса: $p(\theta)\propto\sqrt{\det I(\theta)}$, инвариантно к репараметризации. Для Бернулли $p(p)\propto p^{-1/2}(1-p)^{-1/2}$ (Beta$(\tfrac12,\tfrac12)$). Для среднего нормального~--- $p(\mu)\propto 1$. Для $\sigma$~--- $p(\sigma)\propto 1/\sigma$.
\begin{itemize}
    \item Априор, противоречащий механизму генерации данных, смещает вывод
          даже при больших $n$ (несобственные априоры требуют проверки нормируемости апостериора).
    \item Достоверный и доверительный интервалы по формулировке
          часто путают. 95\% ДИ~--- свойство процедуры. 95\% достоверный~--- вероятность параметра после наблюдения данных.
    \item Несобственный равномерный на $\mathbb{R}$ для дисперсии не даёт нормируемого
          апостериора в малых выборках.
\end{itemize}

Выбор априора~--- часть постановки. $\mathrm{Beta}(2,2)$ кодирует умеренное предположение о справедливости монетки до данных. $\mathrm{Beta}(1,1)$~--- максимальную неопределённость. Анализ чувствительности повторяет вывод при трёх априорах. Если MAP и HPD радикально различаются при $n=20$, данные ещё не доминируют~--- это нормально сообщать в отчёте, а не подбирать априор.

\section{Связь с регуляризацией}

Для нормального правдоподобия с гауссовым априором на коэффициенты: \[ \log p(\beta\mid X,y)\propto -\frac1{2\sigma^2}\|y-X\beta\|_2^2-\frac1{2\tau^2}\|\beta\|_2^2 \] ~--- ridge ($l_2$). Априор Лапласа $\Rightarrow$ lasso ($l_1$).
Регуляризация Фирта в логистической регрессии~--- байесовский штраф Джеффриса-подобного вида при разделимости классов (линейно разделимые классы).
Интерпретация. $\tau^2$ в ridge~--- дисперсия априора на $\beta$. Чем меньше $\tau$, тем сильнее сжатие коэффициентов к нулю.
Lasso соответствует априору Лапласа с экспоненциальными хвостами, что объясняет разреженность MAP-оценок.
В гл.~\ref{ch:regression} ridge MAP $\Leftrightarrow$ $l_2$-штраф, lasso MAP $\Leftrightarrow$ $l_1$. Байесовский вывод делает штраф
интерпретируемым как $\log p(\beta)$. При мультиколлинеарности гауссов априор стабилизирует $\hat\beta_{\mathrm{MAP}}$ так же, как ridge стабилизирует МНК~--- различие в интерпретации (распределение и частотная процедура), сходство в формуле.

\section{Достоверные интервалы}

\begin{definition}
$(1-\alpha)$-достоверный интервал $[a,b]$ удовлетворяет $P(\theta\in[a,b]\mid X)=1-\alpha$.
\end{definition}

С равными хвостами: квантили $\alpha/2$ и $1-\alpha/2$ апостериора.
HPD (наибольшая апостериорная плотность) даёт кратчайший интервал с массой $1-\alpha$.
Для $\mathrm{Beta}(8,4)$ HPD смещён относительно интервала с равными хвостами при асимметрии. При $\alpha+\beta+n$ большом различие мало.
В отчёте по практической задаче (гл.~\ref{ch:labs}) для вероятностных моделей предпочтительно указывать и точечную оценку (среднее/MAP), и HPD~--- особенно при асимметричном апостериоре и малых $n$.
Чувствительность к априору. При $n=10$, $\bar X=0{,}17$, $X_i\sim\mathcal{N}(0,1)$ частотный 95\% ДИ для $\mu$~--- $[-0{,}45,\,0{,}78]$, а при сильном априоре $\mu\sim\mathcal{N}(0,0{,}01)$ достоверный~--- $[-0{,}002,\,0{,}03]$. При $n=10^5$ оба сужаются к $[0{,}164,\,0{,}176]$, но априор всё ещё слегка смещает центр.
Анализ чувствительности повторяет вывод при нескольких разумных априорах (информативный, слабый, Джеффриса) и сравнивает MAP, достоверный интервал и апостериорное прогнозное распределение. Существенный разброс при фиксированном $n$~--- сигнал, что данные не доминируют над априором.
Практически достоверный интервал строят из сэмплов MCMC или аналитически для сопряжённых семейств. С равными хвостами проще, HPD короче при асимметрии. В отчёте указывают тип интервала. 95\% достоверный с равными хвостами и 95\% HPD могут слегка отличаться при $\mathrm{Beta}(8,4)$.

\section{Аппроксимативный вывод: MCMC}

Когда $p(\theta\mid X)$ не вычисляется в замкнутой форме, сэмплируют из апостериора.

Metropolis--Hastings:.\begin{enumerate}
    \item Предложение $z'\sim q(z'\mid z^{(t)})$.
    \item Принять с вероятностью
          $A=\min\!\left(1,\dfrac{p(z')\,q(z^{(t)}\mid z')}{p(z^{(t)})\,q(z'\mid z^{(t)})}\right)$.
    \item Если принято, $z^{(t+1)}=z'$; иначе $z^{(t+1)}=z^{(t)}$.
\end{enumerate}
Достаточно знать $p$ с точностью до нормировки. Цепь имеет стационарное распределение $p(z\mid X)$ при симметричном $q$ (случайного блуждания Метрополиса).
MCMC~--- частотный способ получить байесовские интегралы. Среднее по сэмплам аппроксимирует $\mathbb{E}(\theta\mid X)$, квантили~--- достоверные интервалы.
Сэмплы коррелированы. Для оценок берут выход на стационарность и прореживание (каждый $k$-й сэмпл). Диагностика включает графики траекторий. $\hat R>1$ указывает на плохую сходимость цепей.
Доля принятия и длина выхода на стационарность зависят от формы апостериора. Типичная последовательность шагов. Задать априор и правдоподобие $\to$ запустить цепь $\to$ проверить траекторию и $\hat R$ $\to$ построить достоверный интервал и прогнозное распределение $\to$ при необходимости сравнить модели через BF или BIC.
Без диагностики сходимости байесовский интервал~--- иллюзия точности.

\section{Выбор моделей}

Пусть $\mathcal{M}_k$~--- модель с параметром $\theta_k$.
Фактор Байеса:.\[ BF_{10}=\frac{p(X\mid\mathcal{M}_1)}{p(X\mid\mathcal{M}_0)}= \frac{\int p(X\mid\theta_1,\mathcal{M}_1)p(\theta_1\mid\mathcal{M}_1)\,\mathrm{d}\theta_1} {\int p(X\mid\theta_0,\mathcal{M}_0)p(\theta_0\mid\mathcal{M}_0)\,\mathrm{d}\theta_0}.
          \]

\begin{center}
\small
\begin{tabular}{cl}
\toprule
$BF_{10}$ & Сила доказательств (Джеффриса) \\
\midrule
$<1$ & доказательства в пользу $\mathcal{M}_0$ \\
$1$--$3$ & слабые \\
$3$--$10$ & умеренные \\
$10$--$30$ & сильные \\
$30$--$100$ & очень сильные \\
$>100$ & решительные \\
\bottomrule
\end{tabular}
\end{center}

Evidence $p(X\mid\mathcal{M})=\int p(X\mid\theta)p(\theta\mid\mathcal{M})\,\mathrm{d}\theta$ штрафует сложные модели (Occam factor).
Интуиция. Модель с большим числом параметров может подогнать данные лучше, но размазанный по $\theta$ априор снижает среднее правдоподобие после интегрирования~--- тот же дух, что штраф $k$ или $k\ln n$ в AIC/BIC (гл.~\ref{ch:regression}).
AIC $=-2\log L+2k$~--- асимптотическая аппроксимация для слабых априоров. BIC $=-2\log L+k\ln n$~--- ближе к фактор Байеса при априор с единичной информацией.
MDL: $\mathrm{MDL}(\mathbf{f},\mathfrak{D})=L(\mathbf{f})+L(\mathfrak{D}\mid\mathbf{f})$. При больших $n$ связан с BIC.
Savage--Dickey (вложенные модели). Если $\mathcal{M}_0$ получается из $\mathcal{M}_1$ при $\theta=\theta_0$, то \[ BF_{10}=\frac{p(\theta_0\mid\mathcal{M}_1,X)}{p(\theta_0\mid\mathcal{M}_1)}.
\] Нужна только апостериорная плотность под $\mathcal{M}_1$ в точке $\theta_0$ и априор~--- удобно для сравнения эффекта / отсутствия эффекта в регрессии (гл.~\ref{ch:regression}). $\mathcal{M}_0$~--- модель без предиктора, $\mathcal{M}_1$~--- с ним. $\theta_0=0$ для коэффициента.
Аппроксимация Лапласа: $\log p(X\mid\mathcal{M})\approx \log p(X\mid\hat\theta)\,p(\hat\theta\mid\mathcal{M}) -\frac{k}{2}\log(2\pi)+\frac12\log\det H^{-1}$, где $H$~--- гессиан $-\log p(\theta\mid X)$ в моде.
Для $g$-априор на $\beta$ в линейной регрессии маргинальное правдоподобие выписывается аналитически.
Вычисление маргинального правдоподобия многомерно. Оценка гармонического среднего MCMC может иметь бесконечную дисперсию. Для вложенных моделей предпочтительны Savage--Dickey или мостовая выборка. BF чувствителен к априору на общих параметрах вложенных моделей (парадокс Линдли при неинформативном априоре на $\theta$).
Линейная регрессия с $g$-априор даёт $p(X\mid\mathcal{M})$ аналитически. Нелинейные модели~--- аппроксимация Лапласа в моде $\hat\theta$ с гессианом $H$. Для не вложенных моделей~--- мостовая выборка или BIC при описании в отчёте. Гармоническое среднее из MCMC для маргинального правдоподобия ненадёжно.
Пример выбора модели в регрессии. $\mathcal{M}_0$ без $x_j$, $\mathcal{M}_1$ с $x_j$. Savage--Dickey даёт $BF_{10}$ через апостериор $\hat\beta_j$ при $\mathcal{M}_1$ в точке 0. Согласование с частичным $F$ из гл.~\ref{ch:regression} при слабом априоре и большом $n$. Расхождение при информативном априоре~--- повод для анализа чувствительности.

\section{Evidence и принцип MDL}

\begin{center}
\small
\begin{tabular}{p{0.42\linewidth}|p{0.42\linewidth}}
\toprule
Evidence $p(X\mid\mathcal{M})=\int p(X\mid\theta)\,p(\theta\mid\mathcal{M})\,\mathrm{d}\theta$
& MDL $\mathrm{MDL}(\mathbf{f},\mathfrak{D})=L(\mathbf{f})+L(\mathfrak{D}\mid\mathbf{f})$ \\
\midrule
Использует априорные знания через $p(\theta\mid\mathcal{M})$
& Минимизирует длину описания данных в битах. Априор не обязателен \\
\midrule
Опирается на вероятностную модель порождения выборки
& Сравнивает модели по коду без явной генеративной модели \\
\midrule
фактор Байеса, Savage--Dickey, Laplace
& Связь с BIC при больших $n$: $\mathrm{BIC}\approx -2L+\log n\,(k{+}1)$ \\
\bottomrule
\end{tabular}
\end{center}

При большом $n$ MDL разлагается на три слагаемых (длины в битах): \[ \mathrm{MDL}(\mathbf{f},\mathfrak{D}) \approx \underbrace{L(\mathbf{f})}_{\text{сложность модели}} + \underbrace{L(\mathbf{w}^*\mid\mathbf{f})}_{\text{код параметров}} + \underbrace{L(\mathfrak{D}\mid\mathbf{w}^*,\mathbf{f})}_{\text{остаток / шум}}, \] где $\mathbf{w}^*$~--- МП- или MAP-параметры при фиксированной структуре $\mathbf{f}$.
Первое слагаемое штрафует богатые модели (Occam). Третье~--- fit данных.
Evidence $p(X\mid\mathcal{M})$ интегрирует по $\theta$ и даёт тот же дух штрафа за сложность (Occam factor), но явно через априор.

Колмогоровская сложность (длина минимального программы для $\mathfrak{D}$) невычислима. MDL~--- вычислимый практический принцип. Теорема инвариантности даёт близость кодов на разных языках только асимптотически при больших $n$.

\section{Схема байесовского анализа}

\begin{enumerate}
    \item Задать правдоподобие $p(X\mid\theta)$ и осмысленный априор $p(\theta)$.
    \item Аналитический апостериор, сетку, вариационный вывод или MCMC (MH, HMC).
    \item Проверить сходимость ($\hat R$, графики траекторий).
    \item Сообщить среднее/медиана, 95\% достоверный, HPD. Апостериорное прогнозное распределение
          для прогноза.
    \item При нескольких $\mathcal{M}_k$~--- фактор Байеса (Savage--Dickey,
          Laplace, BIC) и анализ чувствительности по априору.
\end{enumerate}
Схема повторяет логику практических задач (гл.~\ref{ch:labs}). Сначала явная модель и предпосылки, затем оценивание, диагностика и интерпретация с учётом альтернативных априоров.
Для HMM (гл.~\ref{ch:causality}) EM даёт MAP-подобную точку. Полный байесовский HMM потребовал бы MCMC по $(A,B,\pi)$~--- тема за рамками курса, но концептуально та же.
В отчёте указывают априор (семейство и параметры), правдоподобие, апостериорное распределение (аналитическое/MCMC), точечную оценку (среднее/MAP), 95\% HPD, апостериорное прогнозное распределение при необходимости. Чувствительность~--- 2--3 априора. BF или BIC между моделями. Не используют байесовское p-value~--- вместо этого хвост апостериорного распределения или BF. При MCMC указывают число сэмплов, выход на стационарность и $\hat R$.

\begin{codefragment}
# Beta-Binomial posterior (7 heads, 3 tails)
import numpy as np
from scipy import stats
flips = ["H"] * 7 + ["T"] * 3
n_heads = sum(1 for x in flips if x == "H")
a0, b0 = 1, 1  # uniform prior
a_post = a0 + n_heads
b_post = b0 + len(flips) - n_heads
ci = stats.beta.ppf([0.025, 0.975], a_post, b_post)
pred = a_post / (a_post + b_post)  # P(next H | data)
\end{codefragment}

Пример: монетка, регрессия, выбор модели (1)~7H/3T, $\mathrm{Beta}(1,1)$ априор. Апостериорное распределение $\mathrm{Beta}(8,4)$. 95\% HPD для $p$ не совпадает с частотным интервалом Уилсона~--- разная интерпретация (гл.~\ref{ch:basics}).
(2)~Линейная регрессия с $g$-априор на $\beta$. MAP$=$ridge. При $H_0\colon\beta_j=0$ Savage--Dickey даёт $BF_{10}$ без интеграла по всему пространству.
(3)~Сравнение моделей с 5 и 12 предикторами. BIC и Laplace-аппроксимация маргинального правдоподобия штрафуют сложность. Лучший $R^2$ на обучении может проигрывать по маргинальному правдоподобию на тесте.

Анализ чувствительности по априору. $\mathrm{Beta}(1,1)$ и $\mathrm{Beta}(2,2)$ при $n=10$ заметно сдвигает HPD. При $n=200$ апостериоры почти совпадают~--- данные подавляют слабый априор. MCMC для логистической регрессии со слабо информативным априором на $\beta$ даёт достоверные интервалы для OR, сопоставимые с Уолда, но устойчивее при разделимости классов.

На практике команды смешивают подходы. MLE/МНК для точечной оценки, бутстреп-ДИ для частотной неопределённости, слабый априор для регуляризации, BIC для выбора модели~--- всё это байесовски мотивировано, но не полный апостериорный отчёт. Полный байесовский отчёт добавляет явный априор, апостериорное распределение (аналитическое или MCMC), HPD, апостериорное прогнозное распределение и BF между моделями. Важно понимать, какой уровень байесовскости использован. Путаница между достоверным интервалом и доверительным~--- частая ошибка (см. гл.~\ref{ch:basics}).
При $n\gg 1$ априор подавляются данными, и подходы сходятся. При малых $n$ различие принципиально.
Слабые априоры и большие $n$ дают апостериор, близкий к правдоподобию, и достоверный интервал $\approx$ доверительный. Малые $n$, иерархические модели, явный учёт неопределённости параметра, сравнение моделей через маргинальное правдоподобие, регуляризация с интерпретацией априора~--- сильные мотивации. Несобственные априоры и MCMC без диагностики~--- зона повышенного риска.
Байесовский вывод~--- не отдельная статистика, а явное объединение правдоподобия, априора и прогнозного распределения. Связь с ridge, AIC/BIC и EM (гл.~\ref{ch:causality}) замыкает курс. В отчёте всегда указывают тип интервала (достоверный или доверительный) и проверку чувствительности к априору при малых $n$.

\section{Практика и семинар}
\label{sec:practice-ch8}

Семинар~13 проходит полный байесовский цикл на синтетических монетках с конвейера с визуализацией сопряжённого обновления.

Синтетические монетки. Постановка. Монетки с конвейера. Доля честных с $p=0{,}5$ и смещённых.
Для одной монетки: $X_i\mid p\sim\mathrm{Bern}(p)$ i.i.d., априор $p\sim\mathrm{Beta}(\alpha,\beta)$.

Шаги вычисления:.\begin{enumerate}
    \item Априор. Задайте $(\alpha,\beta)$. $\mathrm{Beta}(1,1)$ --- неинформативный.
          $\mathrm{Beta}(2,2)$ --- умеренное смещение к $0{,}5$.
    \item Данные. $n$ подбрасываний, $k$ орлов. Правдоподобие
          $p(X\mid p)\propto p^k(1-p)^{n-k}$.
    \item Апостериор.
          $p\mid X\sim\mathrm{Beta}(\alpha+k,\,\beta+n-k)$.
    \item Сводки: $\mathbb{E}(p\mid X)=\dfrac{\alpha+k}{\alpha+\beta+n}$.
          MAP при $\alpha,\beta>1$: $\hat p_{\mathrm{MAP}}=\dfrac{\alpha+k-1}{\alpha+\beta+n-2}$.
    \item Интервалы: 95\% достоверный интервал с равными хвостами из квантилей Beta.
          HPD короче при асимметрии (особенно при малых $n$).
    \item Апостериорное прогнозное распределение:
          $P(X_{\mathrm{new}}=1\mid X)=\dfrac{\alpha+k}{\alpha+\beta+n}$.
    \item Анализ чувствительности: повторите шаги 3--6 для трёх априоров. Сравните HPD и MAP.
\end{enumerate}

Анимация обновления (sem13). При последовательном поступлении орлов плотность $\mathrm{Beta}(\alpha+k,\beta+n-k)$ сдвигается вправо. При $P=0{,}5$ генератора и $\mathrm{Beta}(2,2)$ апостериор сходится к $0{,}5$ при $n\to\infty$.

\begin{codefragment}
# Posterior Beta-Binomial (sem13)
import numpy as np
from scipy import stats
k, n = 7, 10          # observed heads and n
a0, b0 = 1, 1         # Beta(1,1) prior
a_post, b_post = a0 + k, b0 + n - k
mean_p = a_post / (a_post + b_post)
hpd = stats.beta.interval(0.95, a_post, b_post)
pred_next = mean_p      # P(next head | data)
\end{codefragment}

Нормальная модель (дополнение sem13). При известной $\sigma^2$ сопряжённый априор $\mu\sim\mathcal{N}(\mu_0,\sigma_0^2)$ даёт $\mu\mid X\sim\mathcal{N}(\mu_n,\sigma_n^2)$ с $\mu_n=\bigl(\frac{\mu_0}{\sigma_0^2}+\frac{n\bar x}{\sigma^2}\bigr) \bigl(\frac{1}{\sigma_0^2}+\frac{n}{\sigma^2}\bigr)^{-1}$.
Сравните 95\% достоверный для $\mu$ с частотным $t$-интервалом при том же $\alpha$.

Отчёт: укажите априор, правдоподобие, апостериорное распределение (график), среднее/MAP, 95\% HPD, прогнозное распределение для одного нового наблюдения. Не путайте достоверный интервал с ДИ (гл.~\ref{ch:basics},~\ref{ch:hypothesis}).

\section{Контрольные вопросы по теме}

\begin{enumerate}
    \item Сформулируйте различие между достоверным интервалом и доверительным
          интервалом (см. также гл.~\ref{ch:basics}).
    \item Покажите, что $\mathrm{Beta}(1,1)$ априор для монетки даёт апостериор
          $\mathrm{Beta}(1+k,\,1+n-k)$ при $k$ орлах.
    \item Запишите апостериор NIG для $(\mu,\sigma^2)$ в одной строке
          для $\mu_n$.
    \item Выведите апостериорное прогнозное распределение для Beta--Binomial.
    \item Почему $\log p(\theta\mid X)$ можно трактовать как правдоподобие плюс штраф?
    \item Запишите долю принятия Metropolis--Hastings.
    \item Что такое фактор Байеса и как он связан с маргинальным правдоподобием?
    \item Как применить Savage--Dickey для $H_0\colon\beta=0$ в регрессии?
    \item Интерпретируйте строку $BF_{10}\in(3,10)$ по шкале Джеффриса.
    \item Когда BIC и фактор Байеса дают согласованный выбор модели?
    \item Зачем нужно распределение Джеффриса, если есть равномерный априор?
    \item Как интерпретировать апостериорное среднее $p=0{,}7$ против MAP $p=0{,}68$
          при одном и том же априоре?
    \item Приведите пример, когда сильный априор и малый $n$ радикально
          меняют достоверный интервал.
    \item Чем VI отличается от MCMC по ошибке аппроксимации?
\end{enumerate}
